\section{Artifact}
\label{sec:artifact}

Everything needed to repeat this study is in the \texttt{study/} directory of the
\texttt{storage-study} branch of the fork~\cite{forkbranch}.
\begin{itemize}
\item \texttt{study/lib/fpstudy/}: the single-threaded rewrite (\texttt{fastpath.py}),
  the ladder runner, loop-device, dm-delay and NBD helpers, corpus and trace
  generators, and the statistics.
\item \texttt{study/experiments/}: one driver per experiment (E0--E8), each writing one
  JSON line per run with the commit hashes of the code under test, the kernel
  release and the full configuration.
\item \texttt{study/data/}: the raw results behind every number in this paper,
  including the captured traces, and \texttt{EXCLUDE.json}, which lists the rows we
  dropped and why.
\item \texttt{study/analysis/report.py}: regenerates every table, every figure and
  \texttt{results/numbers.json}, which maps each number quoted in the text to the
  file it came from; \texttt{study/\allowbreak analysis/\allowbreak perturb.py} checks the raw data for
  windows of external load (\S\ref{sec:method}).
\end{itemize}
The code under test is exported unmodified from git (upstream 796df86, fork
c162012). Each experiment starts with one command,
\texttt{study/bin/run\_exp.sh <driver> [options]}, which copies the code, prepares BCC's
kernel headers, runs the driver in a private mount namespace with the file system
under test at \texttt{/mnt}, and copies the results back. \texttt{study/README.md}
lists the packages and the exact invocations; \texttt{study/bin/run\_batch1.sh},
\texttt{run\_batch2.sh} and \texttt{run\_batch3.sh} run every local experiment after
E0 and E1, one at a time, in about ten hours on our machine; \texttt{run\_batch4.sh} and
\texttt{run\_batch5.sh} run the exploratory ones of \S\ref{sec:explore} in under an
hour, and \texttt{remote\_batch\_e6.sh} the real-disk subset in about 25 minutes. Their
logs are in \texttt{study/data/logs/}.
